\section{Feedback: when the prediction changes what is predicted}\label{sec:feedback}

So far the block that follows the context was fixed before the prediction was used. In H\"anni's first
suggestion the actions guided by the prediction can show up in the sensor. We model one round of this by a
\emph{response map} $R\colon\B^m\to\B^m$: the predictor is given the context $x$ and returns $Y$; the world
acts on $Y$; the next $m$ bits of the record are $R(Y)$. The earlier sections are the case of constant $R$.
Let $\Fix(R)=\{y:R(y)=y\}$, the self-fulfilling predictions.

\begin{theorem}[One round with feedback]\label{thm:feedback}
Let $Y\sim\pi_m(\cdot\mid x)$ and let the record continue with $R(Y)$.
\begin{enumerate}[label=(\alph*),leftmargin=2em,itemsep=2pt,topsep=2pt]
\item The prediction comes true with probability $\pi_m(\Fix(R)\mid x)$.
\item \emph{Lock-in.} If $Y=y\in\Fix(R)$, then for every $d\ne y$ and every later context $x'\sqsupseteq xy$,
$E(x')\cap E(xd)=\emptyset$: the programs that predicted $d$ have posterior weight $0$ from then on, and the
posterior is the prior restricted to $E(xy)$. If $d$ is itself a fixed point, the block the world would have
shown had the predictor said $d$, this happens with probability $\pi_m(\Fix(R)\setminus\{d\}\mid x)$.
\item \emph{Obedient world.} Suppose $R$ is the identity in every round, so that the record after $x$ is
$Y_1Y_2\cdots$ with $Y_{i+1}\sim\pi_m(\cdot\mid xY_1\cdots Y_i)$, for as long as $Z_m$ is positive at the
context reached. Then for every $u\in\B^{jm}$ the record
continues with $u$ with probability $\prod_{i\le j}\pi_m(u_i\mid xu_{<i})\ge\xi(xu)/\xi(x)$, where $u_i$ is the
$i$-th block of $u$; and every finite part of the realised record has a $(T,L)$-program.
\item \emph{A default that is too expensive is abandoned.} In (c), if $\xi(xd)=0$ for a block $d$, the record
continues with $d$ with probability $0$.
\end{enumerate}
\end{theorem}
\begin{proof}
(a) is the definition. (b) $x'$ extends $xy$, which is incomparable with $xd$; use \cref{lem:basic}(i). (c) The
product formula is the definition of the process. Each factor is at least $\xi(xu_{\le i})/\xi(xu_{<i})$, and
these ratios multiply to $\xi(xu)/\xi(x)$. A realised prefix has positive probability, hence positive $\xi$,
hence a program, by the proof of \cref{lem:basic}(iii). (d) $\pi_m(d\mid x)=0$.
\end{proof}

Three remarks. First, with feedback the mass off the default is no longer a mass of errors. In the setting of
\cref{thm:budget} a deviating sample is contradicted by the next data and its programs lose their weight; here
a deviating sample that is a fixed point is confirmed, and it is the default that loses its weight, for good.
The probability of this in one round is still a posterior mass, so the floors and caps of
\cref{sec:floor,sec:limits} apply to it unchanged: it is at least $2^{-c}\rho$ for a payload that is a fixed
point and that a wrapper can print, and it is small on average over cut points, by \cref{thm:budget}, when the
default is within the bounds. Feedback changes what a deviation \emph{does}, not how likely it is.

Second, (c) and (d) are the precise form of the ``dark room''. An obedient world driven by this predictor
does not seek darkness. It follows a sample from the predictor's own bounded prior conditioned on the past.
That sample stays on the default future with probability at least $\xi(xd_1d_2\cdots)/\xi(x)$, which is
close to $1$ when the default is cheap and much data has been seen. But if the default future needs more time
than the schedule allows, it has probability $0$ and the world is moved, at the first round whose default
block is off schedule, onto a future
that some program of length at most $L$ prints on schedule: constant, periodic, or otherwise cheap. The pull
is toward what is cheap to compute, and its strength is set by the time bound.

Third, whether a prediction can fulfil itself is a matter of timing in the world. If the bits being predicted
are fixed before the action can reach the sensor, $R$ is constant and nothing is self-fulfilling. The
counterfactual oracle of~\cite{armstrong17} enforces this by design; \cite{oesterheld23,hubinger23} treat
predictors that are scored on outcomes they influence.

\begin{theorem}[Contrarian world]\label{thm:contrarian}
Let the predictor return the more probable next bit, $b(x)=1$ if $\xi(x1)>\xi(x0)$ and $b(x)=0$ otherwise, and
let the record satisfy $x_{k+1}=1-b(x_{1:k})$ for all $k\ge n$. Then $\xi(x_{1:n+j})\le2^{-j}\xi(x_{1:n})$ for
all $j$, and if $L<\infty$ then $\xi(x_{1:n+j})=0$ for all $j>L$.
\end{theorem}
\begin{proof}
The bit $1-b(x)$ has $\xi(x(1-b(x)))\le\xi(xb(x))$, and the two sum to at most $\xi(x)$, so
$\xi(x(1-b(x)))\le\xi(x)/2$. If $\xi(x_{1:n+j})>0$ then $2^{-L}\le\xi(x_{1:n+j})\le2^{-j}$ by
\cref{lem:basic}(iii), so $j\le L$.
\end{proof}

With finite bounds this world is computable, and its description is that of the predictor plus a constant; it
is not printed on schedule by any program of length at most $L$, because it contains the predictor, which runs
all of them. This is Putnam's diagonal argument~\cite{putnam63,legg06} placed inside the world; see
also~\cite{wyeth25}, where the unbounded mixture fails to learn a sequence that contains adversarially chosen
actions. After at most $L+1$ rounds
the predictor has no hypothesis left and its output is whatever its implementation does with an empty
posterior. A world that contains a bounded predictor and reacts against it is outside every guarantee above.
